\section{Introduction to Explainable AI}
\subsection{Taxonomy}
By Interpretability Type
\begin{itemize}
	\item Intrinsic Interpretability
	\item Post-hoc Interpretability
\end{itemize}
By Result Type
\begin{itemize}
	\item Feature summary statistics
	\item Feature summary visualizations
	\item Model internals (weights, tree structure)
	\item Data Point base (counterfactuals, prototypes)
	\item Intrinsic surrogates
\end{itemize}
\subsection{Global vs Local interpretability}
\paragraph{Global Interpretability}
How well the trainied model makes predictions overall.
Which features are important for predictions overall.
Feature interactions.
\paragraph{Local Interpretability}
Why the model made a particular prediction for one instance
Locally a complex model may behave in a simple way.

\subsection{Making explanations}
When making explanations, we need to keep the human factor in mind.
They should be
\begin{itemize}
	\item Contrastive: Provide “why this, not that?” with a clear reference point.
	\item Selected: Limit to the top 1–3 drivers.
	\item Social: Match audience knowledge and context.
	\item Abnormality: Highlight unusual but influential factors.
	\item Fidelity: Stay faithful to model behavior (and reality when possible).
	\item Beliefs: Acknowledge/prioritize known priors or add constraints if needed.
	\item Generality: Indicate coverage/support when relevant.
\end{itemize}
